\begin{abstract}
We compare a master's thesis reward with a replacement within a shared
cellular UAV control system, V2, using the same Barcelona ray tracing dataset
as the original thesis. The original reward arm, V1.5, shares V2's dynamics,
observations, navigation prior, safety filter, PPO settings and strict sampled
connectivity rules. Five training seeds and 400 fresh shared routes support
the initial comparison. Full V2 achieves \FullStd\% standard and \FullLong\% longer joint
success, versus \OriginalStd\% and \OriginalLong\% for V1.5. The primary
standard difference is \PrimaryDiff\ percentage points, with a 95\% crossed
seed and route bootstrap interval \PrimaryCI. \PrimaryAbstract{}
Full trades fewer handovers for higher delay. Every successful mission meets
the sampled RSS and buffer constraints. These results concern our declared
simulator; the original implementation remains unavailable.
A subsequent validation selected logarithmic delay reward was tested on 1,000 new shared routes with five seeds. Its standard completion difference versus full V2 is -1.440 points [-3.160, 0.320]. The fresh standard comparison does not establish improved completion over full V2.

A separate fixed weight lookahead supervisor has a fresh standard completion difference of +2.000 points [0.320, 3.681] versus full V2. The primary fresh comparison supports improved standard mission completion.

\end{abstract}
\begin{IEEEkeywords}
Cellular UAV, connectivity constraints, reward design, handover, trajectory
control, resource allocation, reinforcement learning, reproducibility.
\end{IEEEkeywords}

\section{Introduction}
A master's thesis studies joint trajectory, cell association and resource
control using Barcelona ray tracing data\cite[pp.~6--7, Sec.~3.1]{thesis}.
Its objective weights delay, interference and handovers\cite[p.~10, Eq.~(11)]{thesis},
with minimum serving RSS and bounded queue constraints\cite[p.~11, Eq.~(12),
C2 and C4]{thesis}. The thesis reports wandering and prioritization of
connectivity over destination completion\cite[p.~19, Sec.~7.2; p.~21, Sec.~8]{thesis}.
Positive reciprocal radio utilities and a fixed approach bonus
\cite[p.~12, Eqs.~(13)--(14)]{thesis} motivate an empirical reward audit,
but do not establish the cause of behavior in unavailable original code.

The central question is whether replacing that reward improves joint mission
success when every other component is the same V2 system. V1.5 denotes the
written original reward on V2, created as a new controlled intervention.
Both arms retain V2's strict sampled connectivity rules, arrival termination,
navigation prior and prospective safety filter. We do not restore the original
continued training after arrival\cite[p.~12, Sec.~5.2]{thesis}. The experiment
therefore measures the reward package conditional on V2, not the complete
original agent or individual reward terms in isolation.

We contribute a transparent reward comparison, fresh matched tests, all training
seeds, explicit failures and reproducible artifacts. PPO\cite{ppo} is an
established algorithm; we claim no new learning method. Thesis citations give
printed pages; add two for the PDF viewer page. We use the same thesis dataset;
the original simulator and trained policy remain unavailable.
We also recover source motivated metrics and test a separately declared
lookahead supervisor using unchanged full V2 weights on a third fresh route set.
